%%%%%%%%%%%%%%%%%%%%%%%%%%%%%%%%%%%%%%%%%
%  My documentation report
%  Objective: Explain what I did and how, in order to help someone continue with the investigation
%
% Important note:
% Chapter heading images should have a 2:1 width:height ratio,
% e.g. 920px width and 460px height.
%
% The images can be found anywhere, usually on sky surveys websites or the
% Astronomy Picture of the day archive http://apod.nasa.gov/apod/archivepix.html
%
% The original template (the Legrand Orange Book Template) can be found here --> http://www.latextemplates.com/template/the-legrand-orange-book
%
% Original author of the Legrand Orange Book Template:
% Mathias Legrand (legrand.mathias@gmail.com) with modifications by:
% Vel (vel@latextemplates.com)
%
% Original License:
% CC BY-NC-SA 3.0 (http://creativecommons.org/licenses/by-nc-sa/3.0/)
%%%%%%%%%%%%%%%%%%%%%%%%%%%%%%%%%%%%%%%%%
 
%----------------------------------------------------------------------------------------
%	PACKAGES AND OTHER DOCUMENT CONFIGURATIONS
%----------------------------------------------------------------------------------------

\documentclass[11pt,fleqn,english,russian]{report} % book or report

\usepackage[top=3cm,bottom=3cm,left=3.2cm,right=3.2cm,headsep=10pt,letterpaper]{geometry} % Page margins

\usepackage{courier}
\usepackage{xcolor} % Required for specifying colors by name
\definecolor{ocre}{RGB}{25,50,150} % Define the orange color used for highlighting throughout the book

\definecolor{tssteelblue}{RGB}{70,130,180}
\definecolor{tsorange}{RGB}{255,138,88}
\definecolor{tsblue}{RGB}{23,74,117}
\definecolor{tsforestgreen}{RGB}{21,122,81}
\definecolor{tsyellow}{RGB}{255,185,88}
\definecolor{tsgrey}{RGB}{200,200,200}

\definecolor{codegreen}{rgb}{0,0.6,0}
\definecolor{codegray}{rgb}{0.5,0.5,0.5}
\definecolor{codepurple}{rgb}{0.58,0,0.82}
\definecolor{backcolour}{rgb}{0.95,0.95,0.92}

\usepackage{listings}
%Code listing style named "mystyle"
\lstdefinestyle{mystyle}{
	backgroundcolor=\color{tsblue!10},   commentstyle=\color{codegreen},
	keywordstyle=\color{magenta},
	numberstyle=\tiny\color{codegray},
	stringstyle=\color{codepurple},
	basicstyle=\ttfamily\footnotesize,
	breakatwhitespace=false,         
	breaklines=true,                 
	captionpos=b, 
	xleftmargin=.05\linewidth,
	keepspaces=true,                 
	numbers=left,                    
	numbersep=5pt,                  
	showspaces=false,                
	showstringspaces=false,
	showtabs=false,                  
	tabsize=2
}

\usepackage{listings}
\lstset{basicstyle=\footnotesize\ttfamily,breaklines=true, style=mystyle}
\lstset{framextopmargin=50pt,frame=bottomline}

\usepackage{graphicx}

% Font Settings
\usepackage{avant} % Use the Avantgarde font for headings
%\usepackage{times} % Use the Times font for headings
\usepackage{mathptmx} % Use the Adobe Times Roman as the default text font together with math symbols from the Sym­bol, Chancery and Com­puter Modern fonts

\usepackage{microtype} % Slightly tweak font spacing for aesthetics
\usepackage[utf8]{inputenc} % Required for including letters with accents
\usepackage[T2A]{fontenc}

% Bibliography
%\usepackage[style=alphabetic,sorting=nyt,sortcites=true,autopunct=true,babel=hyphen,hyperref=true,abbreviate=false,backref=true,backend=biber]{biblatex}
%\addbibresource{bibliography.bib} % BibTeX bibliography file
%\defbibheading{bibempty}{}

%%%%%%%%%%%%%%%%%%%%%%%%%%%%%%%%%%%%%%%%%
% This is based on the Legrand Orange Book
% Structural Definitions File
%
% The original template (the Legrand Orange Book Template) can be found here --> http://www.latextemplates.com/template/the-legrand-orange-book
%
% Original author of the Legrand Orange Book Template::
% Mathias Legrand (legrand.mathias@gmail.com) with modifications by:
% Vel (vel@latextemplates.com)
%
% Original License:
% CC BY-NC-SA 3.0 (http://creativecommons.org/licenses/by-nc-sa/3.0/)
%
%%%%%%%%%%%%%%%%%%%%%%%%%%%%%%%%%%%%%%%%%
%----------------------------------------------------------------------------------------
%	VARIOUS REQUIRED PACKAGES
%----------------------------------------------------------------------------------------

\usepackage{titlesec} % Allows customization of titles

\usepackage{graphicx} % Required for including pictures
\graphicspath{{Pictures/}} % Specifies the directory where pictures are stored

\usepackage{lipsum} % Inserts dummy text

\usepackage{tikz} % Required for drawing custom shapes

\usepackage[english]{babel} % English language/hyphenation

\usepackage{enumitem} % Customize lists
\setlist{nolistsep} % Reduce spacing between bullet points and numbered lists

\usepackage{booktabs} % Required for nicer horizontal rules in tables

\usepackage{eso-pic} % Required for specifying an image background in the title page

%----------------------------------------------------------------------------------------
%	MAIN TABLE OF CONTENTS
%----------------------------------------------------------------------------------------

\usepackage{titletoc} % Required for manipulating the table of contents

\contentsmargin{0cm} % Removes the default margin
% Chapter text styling
\titlecontents{chapter}[1.25cm] % Indentation
{\addvspace{15pt}\large\sffamily\bfseries} % Spacing and font options for chapters
{\color{ocre!60}\contentslabel[\Large\thecontentslabel]{1.25cm}\color{ocre}} % Chapter number
{}  
{\color{ocre!60}\normalsize\sffamily\bfseries\;\titlerule*[.5pc]{.}\;\thecontentspage} % Page number
% Section text styling
\titlecontents{section}[1.25cm] % Indentation
{\addvspace{5pt}\sffamily\bfseries} % Spacing and font options for sections
{\contentslabel[\thecontentslabel]{1.25cm}} % Section number
{}
{\sffamily\hfill\color{black}\thecontentspage} % Page number
[]
% Subsection text styling
\titlecontents{subsection}[1.25cm] % Indentation
{\addvspace{1pt}\sffamily\small} % Spacing and font options for subsections
{\contentslabel[\thecontentslabel]{1.25cm}} % Subsection number
{}
{\sffamily\;\titlerule*[.5pc]{.}\;\thecontentspage} % Page number
[] 

%----------------------------------------------------------------------------------------
%	MINI TABLE OF CONTENTS IN CHAPTER HEADS
%----------------------------------------------------------------------------------------

% Section text styling
\titlecontents{lsection}[0em] % Indendating
{\footnotesize\sffamily} % Font settings
{}
{}
{}

% Subsection text styling
\titlecontents{lsubsection}[.5em] % Indentation
{\normalfont\footnotesize\sffamily} % Font settings
{}
{}
{}
 
%----------------------------------------------------------------------------------------
%	PAGE HEADERS
%----------------------------------------------------------------------------------------

\usepackage{fancyhdr} % Required for header and footer configuration

\pagestyle{fancy}
\renewcommand{\chaptermark}[1]{\markboth{\sffamily\normalsize\bfseries\chaptername\ \thechapter.\ #1}{}} % Chapter text font settings
\renewcommand{\sectionmark}[1]{\markright{\sffamily\normalsize\thesection\hspace{5pt}#1}{}} % Section text font settings
\fancyhf{} \fancyhead[LE,RO]{\sffamily\normalsize\thepage} % Font setting for the page number in the header
\fancyhead[LO]{\rightmark} % Print the nearest section name on the left side of odd pages
\fancyhead[RE]{\leftmark} % Print the current chapter name on the right side of even pages
\renewcommand{\headrulewidth}{0.5pt} % Width of the rule under the header
\addtolength{\headheight}{2.5pt} % Increase the spacing around the header slightly
\renewcommand{\footrulewidth}{0pt} % Removes the rule in the footer
\fancypagestyle{plain}{\fancyhead{}\renewcommand{\headrulewidth}{0pt}} % Style for when a plain pagestyle is specified

% Removes the header from odd empty pages at the end of chapters
\makeatletter
\renewcommand{\cleardoublepage}{
\clearpage\ifodd\c@page\else
\hbox{}
\vspace*{\fill}
\thispagestyle{empty}
\newpage
\fi}

%----------------------------------------------------------------------------------------
%	THEOREM STYLES
%----------------------------------------------------------------------------------------

\usepackage{amsmath,amsfonts,amssymb,amsthm} % For math equations, theorems, symbols, etc

\newcommand{\intoo}[2]{\mathopen{]}#1\,;#2\mathclose{[}}
\newcommand{\ud}{\mathop{\mathrm{{}d}}\mathopen{}}
\newcommand{\intff}[2]{\mathopen{[}#1\,;#2\mathclose{]}}
\newtheorem{notation}{Notation}[chapter]

%%%%%%%%%%%%%%%%%%%%%%%%%%%%%%%%%%%%%%%%%%%%%%%%%%%%%%%%%%%%%%%%%%%%%%%%%%%
%%%%%%%%%%%%%%%%%%%% dedicated to boxed/framed environements %%%%%%%%%%%%%%
%%%%%%%%%%%%%%%%%%%%%%%%%%%%%%%%%%%%%%%%%%%%%%%%%%%%%%%%%%%%%%%%%%%%%%%%%%%
\newtheoremstyle{ocrenumbox}% % Theorem style name
{0pt}% Space above
{0pt}% Space below
{\normalfont}% % Body font
{}% Indent amount
{\small\bf\sffamily\color{ocre}}% % Theorem head font
{\;}% Punctuation after theorem head
{0.25em}% Space after theorem head
{\small\sffamily\color{ocre}\thmname{#1}\nobreakspace\thmnumber{\@ifnotempty{#1}{}\@upn{#2}}% Theorem text (e.g. Theorem 2.1)
\thmnote{\nobreakspace\the\thm@notefont\sffamily\bfseries\color{black}---\nobreakspace#3.}} % Optional theorem note
\renewcommand{\qedsymbol}{$\blacksquare$}% Optional qed square

\newtheoremstyle{blacknumex}% Theorem style name
{5pt}% Space above
{5pt}% Space below
{\normalfont}% Body font
{} % Indent amount
{\small\bf\sffamily}% Theorem head font
{\;}% Punctuation after theorem head
{0.25em}% Space after theorem head
{\small\sffamily{\tiny\ensuremath{\blacksquare}}\nobreakspace\thmname{#1}\nobreakspace\thmnumber{\@ifnotempty{#1}{}\@upn{#2}}% Theorem text (e.g. Theorem 2.1)
\thmnote{\nobreakspace\the\thm@notefont\sffamily\bfseries---\nobreakspace#3.}}% Optional theorem note

\newtheoremstyle{blacknumbox} % Theorem style name
{0pt}% Space above
{0pt}% Space below
{\normalfont}% Body font
{}% Indent amount
{\small\bf\sffamily}% Theorem head font
{\;}% Punctuation after theorem head
{0.25em}% Space after theorem head
{\small\sffamily\thmname{#1}\nobreakspace\thmnumber{\@ifnotempty{#1}{}\@upn{#2}}% Theorem text (e.g. Theorem 2.1)
\thmnote{\nobreakspace\the\thm@notefont\sffamily\bfseries---\nobreakspace#3.}}% Optional theorem note

%%%%%%%%%%%%%%%%%%%%%%%%%%%%%%%%%%%%%%%%%%%%%%%%%%%%%%%%%%%%%%%%%%%%%%%%%%%
%%%%%%%%%%%%% dedicated to non-boxed/non-framed environements %%%%%%%%%%%%%
%%%%%%%%%%%%%%%%%%%%%%%%%%%%%%%%%%%%%%%%%%%%%%%%%%%%%%%%%%%%%%%%%%%%%%%%%%%
\newtheoremstyle{ocrenum}% % Theorem style name
{5pt}% Space above
{5pt}% Space below
{\normalfont}% % Body font
{}% Indent amount
{\small\bf\sffamily\color{ocre}}% % Theorem head font
{\;}% Punctuation after theorem head
{0.25em}% Space after theorem head
{\small\sffamily\color{ocre}\thmname{#1}\nobreakspace\thmnumber{\@ifnotempty{#1}{}\@upn{#2}}% Theorem text (e.g. Theorem 2.1)
\thmnote{\nobreakspace\the\thm@notefont\sffamily\bfseries\color{black}---\nobreakspace#3.}} % Optional theorem note
\renewcommand{\qedsymbol}{$\blacksquare$}% Optional qed square
\makeatother

% Defines the theorem text style for each type of theorem to one of the three styles above
\newcounter{dummy} 
\numberwithin{dummy}{section}
\theoremstyle{ocrenumbox}
\newtheorem{theoremeT}[dummy]{Theorem}
\newtheorem{problem}{Problem}[chapter]
\newtheorem{exerciseT}{Exercise}[chapter]
\theoremstyle{blacknumex}
\newtheorem{exampleT}{Example}[chapter]
\theoremstyle{blacknumbox}
\newtheorem{vocabulary}{Vocabulary}[chapter]
\newtheorem{definitionT}{Definition}[section]
\newtheorem{corollaryT}[dummy]{Corollary}
\theoremstyle{ocrenum}
\newtheorem{proposition}[dummy]{Proposition}

%----------------------------------------------------------------------------------------
%	DEFINITION OF COLORED BOXES
%----------------------------------------------------------------------------------------

\RequirePackage[framemethod=default]{mdframed} % Required for creating the theorem, definition, exercise and corollary boxes

% Theorem box
\newmdenv[skipabove=7pt,
skipbelow=7pt,
backgroundcolor=black!5,
linecolor=ocre,
innerleftmargin=5pt,
innerrightmargin=5pt,
innertopmargin=5pt,
leftmargin=0cm,
rightmargin=0cm,
innerbottommargin=5pt]{tBox}

% Exercise box	  
\newmdenv[skipabove=7pt,
skipbelow=7pt,
rightline=false,
leftline=true,
topline=false,
bottomline=false,
backgroundcolor=ocre!10,
linecolor=ocre,
innerleftmargin=5pt,
innerrightmargin=5pt,
innertopmargin=5pt,
innerbottommargin=5pt,
leftmargin=0cm,
rightmargin=0cm,
linewidth=4pt]{eBox}	

% Definition box
\newmdenv[skipabove=7pt,
skipbelow=7pt,
rightline=false,
leftline=true,
topline=false,
bottomline=false,
linecolor=ocre,
innerleftmargin=5pt,
innerrightmargin=5pt,
innertopmargin=0pt,
leftmargin=0cm,
rightmargin=0cm,
linewidth=4pt,
innerbottommargin=0pt]{dBox}	

% Corollary box
\newmdenv[skipabove=7pt,
skipbelow=7pt,
rightline=false,
leftline=true,
topline=false,
bottomline=false,
linecolor=gray,
backgroundcolor=black!5,
innerleftmargin=5pt,
innerrightmargin=5pt,
innertopmargin=5pt,
leftmargin=0cm,
rightmargin=0cm,
linewidth=4pt,
innerbottommargin=5pt]{cBox}

% Creates an environment for each type of theorem and assigns it a theorem text style from the "Theorem Styles" section above and a colored box from above
\newenvironment{theorem}{\begin{tBox}\begin{theoremeT}}{\end{theoremeT}\end{tBox}}
\newenvironment{exercise}{\begin{eBox}\begin{exerciseT}}{\hfill{\color{ocre}\tiny\ensuremath{\blacksquare}}\end{exerciseT}\end{eBox}}				  
\newenvironment{definition}{\begin{dBox}\begin{definitionT}}{\end{definitionT}\end{dBox}}	
\newenvironment{example}{\begin{exampleT}}{\hfill{\tiny\ensuremath{\blacksquare}}\end{exampleT}}		
\newenvironment{corollary}{\begin{cBox}\begin{corollaryT}}{\end{corollaryT}\end{cBox}}	

%----------------------------------------------------------------------------------------
%	REMARK ENVIRONMENT
%----------------------------------------------------------------------------------------

\newenvironment{remark}{\par\vspace{10pt}\small % Vertical white space above the remark and smaller font size
\begin{list}{}{
\leftmargin=35pt % Indentation on the left
\rightmargin=25pt}\item\ignorespaces % Indentation on the right
\makebox[-2.5pt]{\begin{tikzpicture}[overlay]
\node[draw=ocre!60,line width=1pt,circle,fill=ocre!25,font=\sffamily\bfseries,inner sep=2pt,outer sep=0pt] at (-15pt,0pt){\textcolor{ocre}{R}};\end{tikzpicture}} % Orange R in a circle
\advance\baselineskip -1pt}{\end{list}\vskip5pt} % Tighter line spacing and white space after remark

%----------------------------------------------------------------------------------------
%	SECTION NUMBERING IN THE MARGIN
%----------------------------------------------------------------------------------------

\makeatletter
\renewcommand{\@seccntformat}[1]{\llap{\textcolor{ocre}{\csname the#1\endcsname}\hspace{1em}}}                    
\renewcommand{\section}{\@startsection{section}{1}{\z@}
{-4ex \@plus -1ex \@minus -.4ex}
{1ex \@plus.2ex }
{\normalfont\large\sffamily\bfseries}}
\renewcommand{\subsection}{\@startsection {subsection}{2}{\z@}
{-3ex \@plus -0.1ex \@minus -.4ex}
{0.5ex \@plus.2ex }
{\normalfont\sffamily\bfseries}}
\renewcommand{\subsubsection}{\@startsection {subsubsection}{3}{\z@}
{-2ex \@plus -0.1ex \@minus -.2ex}
{.2ex \@plus.2ex }
{\normalfont\small\sffamily\bfseries}}                        
\renewcommand\paragraph{\@startsection{paragraph}{4}{\z@}
{-2ex \@plus-.2ex \@minus .2ex}
{.1ex}
{\normalfont\small\sffamily\bfseries}}

%----------------------------------------------------------------------------------------
%	HYPERLINKS IN THE DOCUMENTS
%----------------------------------------------------------------------------------------

% For an unclear reason, the package should be loaded now and not later
\usepackage{hyperref}
\hypersetup{hidelinks,backref=true,pagebackref=true,hyperindex=true,colorlinks=false,breaklinks=true,urlcolor= ocre,bookmarks=true,bookmarksopen=false,pdftitle={Title},pdfauthor={Author}}

%----------------------------------------------------------------------------------------
%	CHAPTER HEADINGS
%----------------------------------------------------------------------------------------

% The set-up below should be (sadly) manually adapted to the overall margin page septup controlled by the geometry package loaded in the main.tex document. It is possible to implement below the dimensions used in the goemetry package (top,bottom,left,right)... TO BE DONE

\newcommand{\thechapterimage}{}
\newcommand{\chapterimage}[1]{\renewcommand{\thechapterimage}{#1}}

% Numbered chapters with mini tableofcontents
\def\thechapter{\arabic{chapter}}
\def\@makechapterhead#1{
\thispagestyle{empty}
{\centering \normalfont\sffamily
\ifnum \c@secnumdepth >\m@ne
\if@mainmatter
\startcontents
\begin{tikzpicture}[remember picture,overlay]
\node at (current page.north west)
{\begin{tikzpicture}[remember picture,overlay]
\node[anchor=north west,inner sep=0pt] at (0,0) {\includegraphics[width=\paperwidth]{\thechapterimage}};
%%%%%%%%%%%%%%%%%%%%%%%%%%%%%%%%%%%%%%%%%%%%%%%%%%%%%%%%%%%%%%%%%%%%%%%%%%%%%%%%%%%%%
% Commenting the 3 lines below removes the small contents box in the chapter heading
%\fill[color=ocre!10!white,opacity=.6] (1cm,0) rectangle (8cm,-7cm);
%\node[anchor=north west] at (1.1cm,.35cm) {\parbox[t][8cm][t]{6.5cm}{\huge\bfseries\flushleft \printcontents{l}{1}{\setcounter{tocdepth}{2}}}};
\draw[anchor=west] (5cm,-9cm) node [rounded corners=20pt,fill=ocre!10!white,text opacity=1,draw=ocre,draw opacity=1,line width=1.5pt,fill opacity=.6,inner sep=12pt]{\huge\sffamily\bfseries\textcolor{black}{\thechapter. #1\strut\makebox[22cm]{}}};
%%%%%%%%%%%%%%%%%%%%%%%%%%%%%%%%%%%%%%%%%%%%%%%%%%%%%%%%%%%%%%%%%%%%%%%%%%%%%%%%%%%%%
\end{tikzpicture}};
\end{tikzpicture}}
\par\vspace*{230\p@}
\fi
\fi}

% Unnumbered chapters without mini tableofcontents (could be added though) 
\def\@makeschapterhead#1{
\thispagestyle{empty}
{\centering \normalfont\sffamily
\ifnum \c@secnumdepth >\m@ne
\if@mainmatter
\begin{tikzpicture}[remember picture,overlay]
\node at (current page.north west)
{\begin{tikzpicture}[remember picture,overlay]
\node[anchor=north west,inner sep=0pt] at (0,0) {\includegraphics[width=\paperwidth]{\thechapterimage}};
\draw[anchor=west] (5cm,-9cm) node [rounded corners=20pt,fill=ocre!10!white,fill opacity=.6,inner sep=12pt,text opacity=1,draw=ocre,draw opacity=1,line width=1.5pt]{\huge\sffamily\bfseries\textcolor{black}{#1\strut\makebox[22cm]{}}};
\end{tikzpicture}};
\end{tikzpicture}}
\par\vspace*{230\p@}
\fi
\fi
}
\makeatother % Insert the commands.tex file which contains the majority of the structure behind the template

\begin{document}
\title{kernel\_slicer: руководство пользователя}

%----------------------------------------------------------------------------------------
%	TITLE PAGE
%----------------------------------------------------------------------------------------

\begingroup
\thispagestyle{empty}
\AddToShipoutPicture*{\put(150,200){\includegraphics[scale=0.25]{../images/logo.png}}} % Image background

\AddToShipoutPicture*{\put(0,-375){\includegraphics[scale=1.0]{back_title.jpg}}} % Image

\AddToShipoutPicture*{\put(0,650){\includegraphics[scale=1.0]{back_title.jpg}}} % Image


\centering
\vspace*{5cm}
\par\normalfont\fontsize{35}{35}\sffamily\selectfont
\textbf{kernel\_slicer}\\
{\LARGE Руководство пользователя}\par % Book title
%\vspace*{1cm}
%{\Huge Фролов В.А.}\par % Author name
\endgroup

%----------------------------------------------------------------------------------------
%	TABLE OF CONTENTS
%----------------------------------------------------------------------------------------

\chapterimage{baikal2.jpg} % Table of contents heading image

\pagestyle{empty} % No headers

\tableofcontents % Print the table of contents itself

%\cleardoublepage % Forces the first chapter to start on an odd page so it's on the right

\pagestyle{fancy} % Print headers again

%----------------------------------------------------------------------------------------
%	CHAPTER 1
%----------------------------------------------------------------------------------------

\chapterimage{baikal2.jpg} % Chapter heading image

\chapter{Введение}

\section{Что такое kernel\_slicer?}\index{what}

kernel\_slicer --- это генератор кода, принимающий на вход класс на С++ c некоторыми вычисленими/алгоритмами, и генерирующий на выходе другой класс, реализующий те же самые вычисления/алгоритмы на GPU при помощи некоторой технологии программирования нижнего уровня. В данный момент такими технологиями являются Vulkan для векторизации вычислений на GPU и ISPC \cite{ispc} для векторизации вычислений на CPU. В листинге \ref{lst:add_vectors} показан пример простейшего класса, реализующего сложения двух векторов.

\begin{lstlisting}[language=C++, caption=сложение двух векторов]
class Test 
{
public:	
  void kernel1D_add(const float* a, const float* b, float* res, int N) {
    for(int i=0;i<N;i++)
      res[i] = a[i] + b[i];
  } 
	
  virtual void AddVec(const float* a   [[size("N")]], 
                      const float* b   [[size("N")]], 
                            float* res [[size("N")]], int N) {
		kernel1D_add(a,b,res,N); 
  } 	
};
\end{lstlisting}\label{lst:add_vectors}

Предполагается, что пользователь такого класса работает с функцией ``AddVec'' в своём коде как обычно. Обратите внимание, что эта функция виртуальная. В сгенерированном классе Test\_GPU, наследуемом от Test, она будет замещена реализацией сложения векторов на GPU. Конструкция ``[[size("N")]]'' определяет размер данных, которые необходимо копировать между CPU и GPU в данном примере.

\section{Предыстория}\index{Histotry}

Последнее десятилетие технологии программирования GPU (как и других массивно-параллельных вычислительных систем) шли очень непростой дорогой. Эта дорога, или, вернее, дороги, появились на поле постоянного противоречия между аппаратным ускорением и кросс-платформенностью. 

Когда разработчик полагается на одного производителя вычислительного оборудования (например, Nvidia), вопрос выбора технологии программирования, как правило, остро не стоит. Можно использовать технологии программирования, которые этот производитель поддерживает наилучшим образом \cite{NVCPP,CUDA,OpenACC}. Однако такой подход возможен далеко не всегда. 

Кросс-платформенная разработка (под разные вычислительные системы, центральные процессоры с различными векторными архитектурами, графические процессоры от разных производителей и т.д.)  высоко-производительного ПО представляет собой сложную проблему, для которой в настоящий момент нет простого решения. И тогда и сейчас единственной технологией программирования, которая позволяет осуществить доступ к широкому спектру аппаратного ускорения в GPU на большом числе различных устройств является Vulkan. Однако с Vulkan всё не так просто (рис. \ref{fig:mem}, \ref{fig:vulkan_complexity}). 

\begin{figure}[h]
\begin{minipage}{0.45\textwidth}
	\centering
	\includegraphics[width=\textwidth]{mem.jpg}
	\caption{Известный мем, пошедший из предложения ``нельзя просто так взять и напасть на Мордор'' в фильме ``Властелин Колец''.}
	\label{fig:mem}
\end{minipage}
\hfill
\begin{minipage}{0.45\textwidth}
	\centering
	\includegraphics[width=\textwidth]{vulkan_is_hard.png}
	\caption{Простые программы писать на Vulkan сложно, а сложные -- очень сложно.}
	\label{fig:vulkan_complexity}
\end{minipage}
\end{figure}

%\begin{figure}[h]
%	\centering
%	\includegraphics[width=0.5\textwidth]{mem.jpg}
%	\caption{Известный мем, пошедший из предложения ``нельзя просто так взять и напасть на Мордор'' в фильме ``Властелин Колец''.}
%	\label{fig:mem}
%\end{figure}

К сожалению, работа с этим API как напрямую, так и посредством обёрток/библиотек является довольно трудоёмкой. Например, простейшая отрисовка треугольника может занимать несколько тысяч строк кода. Как следствие растёт стоимость и время разработки. Эта проблема не столько техническая, сколько фундаментальная: она происходит из того, что в Vulkan сочетаются такие противоречивые вещи как кросс-платформенность и аппаратное ускорение. На одних GPU определённое аппаратное ускорение есть, на других его нет,.. на одних оно реализовано таким образом, на других --- по другому. Изначально kernel\_slicer задумывался как инструмент в борьбе со сложностью при работе с Vulkan. Позже он превратился в самостоятельную технологию программирования.

\section{Мотивация вашей работы с kernel\_slicer }\index{Motivation}

kernel\_slicer не просто позволяет Вам ускорить разработку кода для GPU на Vulkan. Он позволяет делать 4 вещи, которые в других технологиях программирования затруднительно сделать одновременно:

\begin{enumerate}

\item Вы можете сохранить описание своих программ в чистом алгоритимческом виде на обычном С++. Параллельные версии для GPU (Vulkan) и векторных CPU (ISPC) генерируются автоматически. При этом исходная версия может собираться любым существующим компилятором C++. Например, Вы можете также использовать прагмы OpenMP для распараллеливания циклов. Таким образом, kernel\_slicer не делает ваш проект зависимым от себя самого.

\item Процесс разработки позволяет Вам двигаться мелкими шагами, поэтапно и плавно модифицировать код таким образом, чтобы он в итоге стал параллельным и мог быть перенесён на GPU. Это больше всего похоже на обычный рефакторинг.

\item Вы можете осуществить доступ к любому аппаратному ускорению, доступному в Vulkan и ISPC, даже если текущая версия kernel\_slicer его ещё не поддерживает. Для этого предусмотрен менанизм доработки сгенерированного кода. То есть технология рассчитана на это изначально.

\item Если в kernel\_slicer реализована аппаратная поддержка некоторой функциональности, но ваш целевой GPU её не поддерживает, Вы можете заместить неподдерживаемую аппаратную функциональность собственной программной реализацией во время генерации кода.
\end{enumerate}

\begin{remark}Имея опыт работы с различными технологиями (CUDA, OpenCL, DX9-11, OpenGL) и значительную кодовую базу на OpenCL, мы долго не хотели делать свою технологию программирования и рассматривали существующие (Halide, Taichi, HIP и др.). К сожалению, как и в самом OpenCL, всё упиралось в слабую аппаратную поддержку или в жёсткую зависимость от одного производителя. Разработка kernel\_slicer была начата конце 2020-ого года как вынужденная мера по спасению существующих алгоритмических наработок, которые могли потерять актуальность из-за отсутствия аппаратного ускорения. По прошествии 3 лет мы можем сказать, что эта мера сработала. Наши программы более не зависят ни от одного производителя GPU, и вобщем-то, даже от Вулкана они не зависят, т.к. при необходимости мы всегда можем добавить в kernel\_slicer почти любой новый бэк-енд. При этом во время разработки алгоритмической части мы, как правило, вообще забываем про существование GPU.
\end{remark}

%This statement requires citation \cite{book_key}; this one is more specific \cite[122]{article_key}.

%----------------------------------------------------------------------------------------
%	CHAPTER 2
%----------------------------------------------------------------------------------------
\chapterimage{take_off2.jpg}

\chapter{Быстрый старт}

\section{Сборка}

\begin{enumerate}
	
\item Склонируйте на свою машину kernel\_slicer и его сабмодули: \begin{itemize}
	\item git clone https://github.com/Ray-Tracing-Systems/kernel\_slicer.git
	\item git submodule init
	\item git submodule update
\end{itemize}

\item Установите в своей системе llvm и clang 17-ой версии: \begin{itemize}
	\item wget https://apt.llvm.org/llvm.sh
	\item chmod +x llvm.sh
	\item sudo ./llvm.sh 17
	\item sudo apt-get install llvm-17-dev
	\item sudo touch /usr/lib/llvm-17/bin/yaml-bench
	\item sudo apt-get install libclang-17-dev
	\item sudo apt install clang-17
\end{itemize}

\item Осуществите сборку при помощи CMake: 
\begin{itemize}
	\item cd kernel\_slicer
    \item mkdir cmake-build-release \&\& cd cmake-build-release
    \item cmake -DCMAKE\_BUILD\_TYPE=Release .. 
    \item make -j 8
\end{itemize}

\end{enumerate}

\begin{remark}
	kernel\_slicer можно собрать и использовать под ОС Windows и другие операционные системы, поскольку он является обычным плагином для clang. Для этого смотрите в репозитории инструкцию по сборке в составе LLVM. В данный момент мы не поддерживаем документацию и примеры для Windows в актуальном состоянии, поэтому не рекомендуем начинать с этого. 
\end{remark}

\pagebreak
\section{Запуск первого примера}

\underline{Шаг 1}: Склонируйте первый туториал из репозитория:

\begin{itemize}
	\item git clone https://github.com/kernel-slicer-tutorials/tutorial1\_vec\_add
	\item cd tutorial1\_vec\_add
	\item git submodule init
	\item git submodule update
\end{itemize}

\vspace*{5px}
Обратите внимание, что туториал, как и сам репозиторий kernel\_slicer использует три существующих репозитория на github: (1) LiteMath для базовых математических операций, (2) vk-utils как легковесную библиотеку хелперов (функций-помошников) для работы с Vulkan и volk (2) для корректной загрузки указателей на функций Vulkan.

\vspace*{5px}
\noindent\underline{Шаг 2}: Соберите программу (пока что на СPU), проанализируйте её исходные коды и убедитесь что она работает: 
\begin{itemize}
	\item mkdir cmake-build-debug \&\& cd cmake-build-debug
	\item cmake -DCMAKE\_BUILD\_TYPE=Debug ..
	\item make -j 2
	\item ./test\_vadd 
\end{itemize}

\vspace*{5px}
\noindent\underline{Шаг 3}: Запустите kernel\_slicer:
\begin{itemize}
	\item Отредактируйте файл run\_kslicer.sh. Измените значение переменной 'kslicer\_directory' так, чтобы она указывала на правильный путь для вашей системы.
	\item bash run\_kslicer.sh 
\end{itemize}

\vspace*{5px}
\noindent\underline{Шаг 4}: Соберите сгенерированный шейдер:
\begin{itemize}
\item cd shaders\_gpu 	
\item bash build.sh
\end{itemize}

\vspace*{5px}
\noindent\underline{Шаг 5}: Соберите программу (в этот раз уже для GPU) и запустите её:
\begin{itemize}
	\item mkdir cmake-build-debug \&\& cd cmake-build-debug
	\item cmake -DCMAKE\_BUILD\_TYPE=Debug -DUSE\_VULKAN=ON ..
	\item make -j 8
	\item cd ..
	\item ./cmake-build-debug/test\_vadd 
\end{itemize}

Обратите внимание, что запуск производится из основной папки туториала для того чтобы SPIR-V файл корректно загрузился. Вы также можете пользоваться предоставленными конфигами VS Code для сборки, запуска и отладки. Попробуйте зайти внутрь функции 'AddVec' в отладчике.

%\vspace*{5px}
%\noindent\underline{Шаг 6 (опционально)}: Проанализируете сгенерированные файлы.
%\vspace*{2px}
%\noindent Файлы хостового кода:
%\begin{enumerate}
%\item test\_gpu.h -- заголовочный файл для сгенерированного класса
%\item test\_gpu.cpp -- основной код на Vulkan
%\item test\_gpu\_init.cpp -- инициализация Vulkan-объектов
%\item test\_gpu\_ds.cpp -- инициализация дескриптор-сетов
%\item include/Test\_gpu\_ubo.h -- выделенная структура используемых полей класса
%\end{enumerate}

%\vspace*{2px}
%\noindent Файлы шейдеров:
%\begin{enumerate}
%\item shaders\_gpu/build.sh -- последовательность команд для сборки шейдеров	
%\item shaders\_gpu/common\_generated.h -- файл, подключаемый во все шейдеры
%\item shaders\_gpu/kernel1D\_add.comp -- файл шейдера сложения векторов
%\end{enumerate}	

\begin{remark} 
Вы могли заметить, что наш туториал использует кастомный флаг в CMake 'USE\_VULKAN' для опционального подключения Вулкана и задействования GPU версии программы. Если Вы плохо знакомы с CMake, то самое время начать знакомиться с ним. В современном мире разработки ПО это абсолютно необходимо. 
\end{remark}

\chapterimage{baikal2.jpg}
\chapter{Процесс разработки и основы}

Процесс разработки который Вы встретите в туториалах, предполагает, что Вы можете делать несколько различных реализаций ваших алгоритмов на разных технологиях программирования для CPU и GPU. Независимо от того, используете ли Вы kernel\_slicer или нет. 

Кроме того, скорее всего вы захотите реализовывать различные варианты алгоритмов -- более быстрые и простые, либо наоборот, более медленные и качественные. Наконец, входные и выходные данные, как правило, лежат в обычной хостовой памяти (более сложный вариант, когда данные уже в памяти GPU мы рассмотрим в главе \ref{vulkan-direct}). Именно поэтому взаимодействие ваших алгоритмов с внешней средой предполагается реализовывать через простой и очевидный способ --- виртуальные функции с входными и выходными массивами данных. Рассмотрим как это выглядит на примере: 

\section{Концепция на примере тон-маппинга}

Предположим, что Вы хотите написать библиотеку тонирующих операторов (tone mapping). Смысл тон-маппинга довольно прост: на входе у нас есть изображение широкого динамического диапазона (т.н. HDR изображение), которое может содержать произвольные значения яркости в пикселах (float). На выходе у нас т.н. LDR изображение, где все яркости могут быть только в интервале от 0 до 255. Такое изображение уже можно посмотреть на экране монитора. Собственно для этого тонирующие операторы в основном и нужны.

Количество каналов для определённости пусть будет 4. Тогда для каждого пикселя четыре float значения на входе будут превращаться в четыре 8-битных значения на выходе. Это довольно стандартное представление и четыре таких 8-битных значения обычно просто пакуют в один uint32\_t. Делается это, как правило, для того, чтобы при работе с памятью загрузка и сохранение одного пиксела гарантированно превращалась бы в одну инструкцию, загружающую или, соотвественно, сохраняющую понятное количество байт. Не будем нарушать эту традицию (листинг \ref{lst:tonemapping_api2}).

\begin{lstlisting}[language=C++, caption=интерфейс алгоритмов тон-маппинга]
struct IToneMapping
{
  virtual void Run(int width, int height, 
                   const float* in_data,
                   uint32_t* out_data) = 0;
                   
  virtual void Update(...) {}                 
};
\end{lstlisting}\label{lst:tonemapping_api2}

Функция Run будет делать всю работу, но возможно в нашем интерфейсе потребуется объявить дополнительные функции (например Update для задания настроек алгоритма или функции синхронизации, если предполагается, что алгоритм может работать асинхронно с основной программой). Обратите внимание, что входной массив объявлен как const, а выходной не содержит этого спецификатора. Это важно на будущее. Мы считаем in-out массивы плохой практикой с неочевидной семантикой, и kernel\_slicer их не поддерживает (любой не const указатель он воспринимает как out). Если Вы хотите, чтобы один и тот же массив подавался как на вход и на выход, сделайте в вашей функции явно два указателя (один const, другой обычный), и подайте один и тот же указатель как на вход так и на выход явно.  

Но пока что вернёмся к интерфейсу. Мы можем реализовать простой клампинг значений, экспоненциальное отображение или, например, известный алгоритм тон-маппинга Рейнхарда \cite{Reinhard05}. Или что-то более сложное \cite{Mantiuk08}. Для каждой такой реализации Вы скорее всего напишете класс, который будет замещать виртуальную функцию Run. Это является примером нормального и естественного использованием ООП в С++. Очевидно, что в таком случае Вы будете пользоваться этим механизмом и для подключения реализаций на GPU. Давайте считать для определённости, что мы хотим реализовать алгоритм Рейнхарда в классе ReinhardTM (листинг \ref{lst:reinhard}).

\begin{lstlisting}[language=C++, 
	               caption=класс для алгоритма Рейнхарда, 
	               label=lst:reinhard]	
class ReinhardTM : public IToneMapping
{
public:
  void Run(int width, int height, 
           const float* in_data  [[size(w*h*4)]], 
           uint32_t*    out_data [[size(w*h)]]) override;
		   
  void Update(...) override { ... }
   	
protected:
  void kernel1D_findMax(int size, const float*);
  void kernel2D_process(int w, int h, 
                        const float* inpData, 
                        uint32_t*    outData);
  float m_whitePoint;
};
\end{lstlisting}

Далее мы захотим написать классы ReinhardTM\_Vulkan и ReinhardTM\_ISPC, которые будут реализовывать оператор Рейнхарда на GPU и CPU, но уже векторизованным образом, при помощи ISPC. Мы имеем полное право сделать это без участия kernel\_slicer, реализовав соответствующие классы через соответсвующие API программирования GPU самостоятельно. Однако эта работа техническая и достаточно трудоёмкая. Кроме этого, нам придётся повторять её каждый раз, когда CPU вермися алгоритма изменится. Для больших программ такая ``поддержка'' может превратиться в настоящий кошмар из-за большого количества деталей, с которыми нам придётся иметь дело. kernel\_slicer позволяет написать только класс ReinhardTM, а параллельные реализации ReinhardTM\_Vulkan и ReinhardTM\_ISPC будут сгенерированы автоматически. Для того чтобы это работало, мы должны следать нескольким правилам при написании классов.

\section{Основные правила написания класса}

kernel\_slicer является плагином для clang. Поскольку единица компиляции в clang это один текстовый файл (обычно с расширением ``.cpp''), Вам прежде всего нужно грамотно разделить код на файлы таким образом, чтобы не подавать в kernel\_slicer тот код, который не должен им обрабатываться. 

\underline{Правило №1:} Выполняйте реализацию Ваших алгоритмов изолированно от остальной части программы в одном или нескольких классах. Для каждого такого класса лучше всего делать один или несколько ``.cpp'' файлов.

\begin{remark}
Если в вашем классе есть код, который не должен обрабатываться kernel\_slicer-ом, просто поместите его в отдельный файл.
\end{remark}

\underline{Правило №2:} Все ваши алгоритмы должны иметь двух-уровневую организацию. Это функции (т.н. управляющие функции), которые вызывают другие функции (т.н. вычислительные ядра). Код, который должен что-то считать реализуйте внутри функций-ядер. Код, который управляет запуском ядер (например в зависимости от значения параметра запускает один или другой кернел) реализуйте в управляющих функциях (листинг \ref{lst:RunFunc}). Код вычислительных ядер выполняется на GPU, код управляющих функций --- на CPU.

\begin{lstlisting}[language=C++, 
	               caption=реализация управляющей функции Run, 
	               label=lst:RunFunc]	
void ReinhardTM::Run(int width, int height, 
                     const float* in_data, 
                     uint32_t* outData)
{
  kernel1D_finMax(in_data, width*height);
  kernel2D_process(width,height,in_data,outData);
}
\end{lstlisting}

\underline{Правило №3:} Вычислительные ядра должны содержать в имени префикс kernel1D\_, kernel2D\_ или kernel3D\_. Число перед буквой D указывает на то, сколько вложенных циклов нужно параллелить. Остальные циклы и все функции вызываемые из них будут вставлены в код шейдера (листинг \ref{lst:kernel2D_process}). 

\underline{Правило №4:} Количество итераций для параллельных циклов должно: (1) передаваться через аргументы функций-ядер либо (2) быть переменной, являющейся членом класса. Если Вам нужно вычислить это количество каким-то образом, Вы можете сделать это двумя путями.
\begin{enumerate}
\item Вычислить это число в управляюшей функции, как в нашем примере, где вычисляется произведение width*height для ядра kernel1D\_findMax. В этом случае эти вычисления будут произведены на CPU.
\item Вычислить это число в другом вычислительном ядре и сохранить его в переменной, являющейся членом класса. Такая переменная должна непосредственно учавствовать далее в выражении for в качестве количества итераций цикла в тех ядрах, которые предполагается запускать этим способом. Забегая вперёд скажем, что это т.н. непрямой вызов вычислительного ядра (indirect dispatch), когда количество итераций/микро-потоков станет известно лишь в момент выполнения программы на GPU. 
\end{enumerate}

\begin{lstlisting}[language=C++, 
	               caption=двумерное вычислительнрео ядро вызывает функцию reinhard\_ext, 
	               label=lst:kernel2D_process]	
float reinhard_ext(float v, float max_white)
{
  float numerator = v * (1.0f + (v / (max_white * max_white)));
  return numerator / (1.0f + v);
}

void ReinhardTM::kernel2D_process(int width, int height, 
                                  const float* in_data, 
                                  uint32_t* outData)
{
  for(int y=0;y<height;y++) // first par. loop (kernel2D_)
  {
    for(int x=0;x<width;x++) // second par. loop (kernel2D_)
    {
      float r = reinhard_ext(in_data[4*(y*width+x)+0], m_whitePoint);
      float g = reinhard_ext(in_data[4*(y*width+x)+1], m_whitePoint);
      float b = reinhard_ext(in_data[4*(y*width+x)+2], m_whitePoint);
			
      int ir  = (int)std::min(r*255.0f, 255.0f);
      int ig  = (int)std::min(g*255.0f, 255.0f);
      int ib  = (int)std::min(b*255.0f, 255.0f);
			
      outData[y*w+x] = 0xFF000000 | (ib << 16) | (ig << 8) | ir;
    }
  }
}
\end{lstlisting}


\underline{Правило №5:} Вы можете передавать данные из управляющей функцией в ядро только через аргументы ядра. Нельзя, например, записать значение переменной члена класса в управляющей функции и прочитать его в ядре: на GPU прочитается старое значение.

\underline{Правило №6:} Вы можете передавать данные между различными вычислительными ядрами либо (1) через их аргументы (как в нашем примере, где оба вычислительных ядра принимают на вход один и тот же указатель in\_data), либо (2) через члены класса. Члены класса могут иметь типы данных std::vector и LiteImage::Image2D.


\chapterimage{baikal2.jpg}
\section{Измеряем время выполнения и копирования}

Ранее мы опустили некоторые детали, в частности Вы могли заметить, что в тестовом классе присутствуют функции CommitDeviceData и GetExecutionTime, которые никак не использовались. 

\subsection{Запуск второго примера}

\noindent\underline{Шаг 1}: Склонируйте второй туториал (фрактал Мандельброта) из репозитория:

\begin{itemize}
	\item git clone https://github.com/kernel-slicer-tutorials/tutorial2\_...
	\item cd tutorial2\_...
	\item git submodule init
	\item git submodule update
\end{itemize}

\vspace*{5px}
\noindent\underline{Шаги 2--5}: повторите шаги 2--5 из главы быстрый старт для туториала фрактала Мандельброта. 

\vspace*{5px}
\noindent\underline{Шаг 6}: Обратите внимание на механизм измерения времени при помощи функции GetExecutionTime.


\chapterimage{baikal2.jpg}
\chapter{Частые проблемы и пути их решения}

Before continuing, first and most importantly you must select the \emph{raw} data you are going to process and later after you acquire experience with an specific dataset the idea is to expand the algorithms to any kind of dataset. The important things are to learn how to input the data correctly, establish the right \emph{learning parameters} in the selected algorithm and find the best way to visualize your results and interpret them correctly.

\chapterimage{baikal2.jpg} % Chapter heading image
\chapter{Особенности шаблона для трассировки лучей}

I discovered surfing on the internet a cloud computing software that is free, has data mining algorithms embedded, is specifically developed for Astronomy and is programmed by Caltech, University Federico II and the Astronomical Observatory of Capodimonte. The homepage website, \url{http://dame.dsf.unina.it/index.html}. Well, the platform for testing is ready!, now what? I requested and account and the next day they sent me an acceptance with my user name and my password approved.
I introduced myself to the documentation, the available clustering functions, the manuals for every method, the blogs and discovered that the was one method available that could work with data cubes and do its clustering on every pattern (number in the multidimensional matrix) which was exactly what I needed. The name of this method is ESOM (Evolving Self Organizing Maps) and I read its manual, did some foolish test with all my image and ... never got a result ... the experiment ran forever (more than two weeks), when I realised that this wasn't the best way to tackle this problem I started considering only clustering on the independent images and not in the data cube due to the fact that the dimensionality was immense. So, in the end my selected methods have some results but not all, here is where all the work has to be done, analysed and tested again.

\chapterimage{baikal2.jpg} % Chapter heading image
\chapter{Параметры командной строки}


\chapterimage{baikal2.jpg} % Chapter heading image
\chapter{Работа с Vulkan-интерфейсом напрямую}\label{vulkan-direct}

\chapterimage{baikal2.jpg} % Chapter heading image
\chapter{Работа с ISPC}


\begin{thebibliography}{10}
	
	\bibitem{NVCPP} NVIDIA. \textit{NVC++: a C++17 compiler for NVIDIA GPUs and AMD, Intel, OpenPOWER, and Arm CPUs.} // 	
	
	\bibitem{CUDA} NVIDIA, Vingelmann, P. \& Fitzek, F.H.P., 2020. CUDA, release: 10.2.89, Available at: https://developer.nvidia.com/cuda-toolkit.	
	
	\bibitem{OpenACC} ``OpenACC'', 2021 URL: https://www.openacc.org/
	
	\bibitem{ispc} Matt Pharr, William R. Mark. ``spc: A SPMD Compiler for High-Performance CPU Programming''
	
	\bibitem{b10} Jacobsen, Niklas. ``LLVM supported source-to-source translation-Translation from annotated C/C++ to CUDA C/C++.'' MS thesis. University of Oslo, 2016. 
	
	\bibitem{b11} Balogh, Gabor Daniel, et al. ``Op2-clang: A source-to-source translator using clang/llvm libtooling'', IEEE/ACM 5th Workshop on the LLVM Compiler Infrastructure in HPC (LLVM-HPC). IEEE, 2018. 
	
	\bibitem{b12} Yang, Po, et al. ``Improving utility of GPU in accelerating industrial applications with user-centered automatic code translation'', IEEE Transactions on Industrial Informatics 14.4, 2017, 1347-1360. 
	
	\bibitem{b13} Pienaar, Jacques A., Srimat Chakradhar, and Anand Raghunathan. ``Automatic generation of software pipelines for heterogeneous parallel systems'', SC'12: Proceedings of the International Conference on HPC, Networking, Storage and Analysis. IEEE, 2012. 	
	
	\bibitem{b14} Wenzel Jakob, Sebastien Speierer, Nicolas Roussel, and Delio Vicini. 2022. DR.JIT: a just-in-time compiler for differentiable rendering. ACM Trans. Graph. 41, 4, Article 124 (July 2022), 19 pages. https://doi.org/10.1145/3528223.3530099
	
	\bibitem{b16} Bakhtin V., Krukov V. ``DVM-Approach to the Automation of the Development of Parallel Programs for Clusters.'' Programming and Computer Software. 45. 121-132. 10.1134/S0361768819030034. 
	
	\bibitem{b18} Ohberg, Tomas. ``Auto-tuning Hybrid CPU-GPU Execution of Algorithmic Skeletons in SkePU'', MS thesis. Linkoping University, 2018.	
	
	\bibitem{b19} Ernstsson, August, Lu Li, and Christoph Kessler. ``SkePU 2: Flexible and type-safe skeleton programming for heterogeneous parallel systems'', International Journal of Parallel Programming 46.1, 2018, 62-80. 
	
	\bibitem{b20} Steuwer, Michel, Philipp Kegel, and Sergei Gorlatch. ``Skelcl-a portable skeleton library for high-level gpu programming'', 2011 IEEE International Symposium on Parallel and Distributed Processing Workshops and Phd Forum. IEEE, 2011. 
	
	\bibitem{b21} SYCL, cross-platform abstraction layer, 2021. URL: https://www.khronos.org/sycl/	
	
	
	\bibitem{b23} Han, Tianyi David, and Tarek S. Abdelrahman. ``hiCUDA: High-level GPGPU programming'', IEEE Transactions on Parallel and Distributed systems 22.1, 2010, 78-90. 	
	
	\bibitem{b24} Wu, Jingyue, et al. ``gpucc: an open-source GPGPU compiler'', Proceedings of the 2016 International Symposium on Code Generation and Optimization. 2016. 	
	
	\bibitem{b25} Sathre, Paul, Mark Gardner, and Wu-chun Feng. ``On the portability of cpu-accelerated applications via automated source-to-source translation'', Proceedings of the International Conference on High Performance Computing in Asia-Pacific Region. 2019.
	
	\bibitem{b26} ``HIP, C++ Runtime API and Kernel Language'', 2021, URL: https://github.com/ROCm-Developer-Tools/HIP	
	
	\bibitem{b27} ``Vulkan specification, indirect dispatch command'', 2021, URL: https://www.khronos.org/registry/vulkan/specs/1.2-extensions/man/html/vkCmdDispatchIndirect.html
	
	\bibitem{b29} Mammeri, Nadjib, and Ben Juurlink. ``Vcomputebench: A vulkan benchmark suite for gpgpu on mobile and embedded gpus'', IEEE International Symposium on Workload Characterization (IISWC). IEEE, 2018. 
	
	\bibitem{b35} Ragan-Kelley, Jonathan, et al. ``Halide: a language and compiler for optimizing parallelism, locality, and recomputation in image processing pipelines''. ACM Sigplan Notices 48.6, 2013, 519-530. 
	
	\bibitem{b37} Adams, Andrew, et al. ``Learning to optimize halide with tree search and random programs'', ACM Trans. Graph. (TOG)38.4,2019:1-12. 
	
	\bibitem{b40} Haidl, Michael, and Sergei Gorlatch, ``PACXX: Towards a unified programming model for programming accelerators using C++14'', LLVM Compiler Infrastructure in HPC. IEEE, 2014.
	
	\bibitem{b41} Haidl, Michael, et al. ``Pacxxv2+ RV: an LLVM-based portable high-performance programming model'', Proceedings of the Fourth Workshop on the LLVM Compiler Infrastructure in HPC, 2017. 
	
	\bibitem{b46} Sean Baxter. ``Circle C++ shaders'', 2021. URL: https://github.com/seanbaxter/shaders  
	
	\bibitem{b48} ``Clang documentation'', 2021. URL: https://clang.llvm.org/docs/LibTooling.html
	
	\bibitem{b49} ``Inja, template engine for modern C++'', 2021. URL: https://github.com/pantor/inja
	
	\bibitem{b42} Sidelnik, Albert, et al, ``Performance portability with the chapel language'', IEEE 26th international parallel and distributed processing symposium. IEEE, 2012. 
	
	\bibitem{b44} Baghdadi, Riyadh, et al. ``Tiramisu: A polyhedral compiler for expressing fast and portable code'', IEEE/ACM International Symposium on Code Generation and Optimization (CGO), IEEE, 2019. 
	
	\bibitem{b30} ``Nvidia OptiX'', 2021, URL: https://developer.nvidia.com/optix 
	
	\bibitem{b31} MS ``DirectML'', 2021: https://github.com/microsoft/DirectML 
	
	\bibitem{b34} Hegarty, James, et al. ``Darkroom: compiling high-level image processing code into hardware pipelines'', ACM Trans.Graph.33.4(2014):144-1. 
	
	\bibitem{b38} Rasch, Ari, Richard Schulze, and Sergei Gorlatch, ``Developing High-Performance, Portable OpenCL Code via Multi-Dimensional Homomorphisms'', Proceedings of the International Workshop on OpenCL, 2019. 
	
	\bibitem{b39} Huang, Tsung-Wei, et al, ``Taskflow: A General-purpose Parallel and Heterogeneous Task Programming System'', IEEE Transactions on Computer-Aided Design of Integrated Circuits and Systems, 2021. 
	
	\bibitem{b45} ``Google clspv. A prototype compiler for a subset of OpenCL C to Vulkan compute shaders.'' 2021. URL: https://github.com/google/clspv
	
	\bibitem{b47} Huang, Shan Shan, David Zook, and Yannis Smaragdakis. ``Morphing: Safely shaping a class in the image of others'', European Conference on Object-Oriented Programming, Springer, Berlin, Heidelberg, 2007. 
	
	\bibitem{b50} Parker, Steven G., et al. ``Optix: a general purpose ray tracing engine'', ACM transactions on graphics (tog) 29.4 (2010): 1-13.
	
	\bibitem{b52} Laine, Samuli, Tero Karras, and Timo Aila. ``Megakernels considered harmful: Wavefront path tracing on GPUs'', Proceedings of the 5th High-Performance Graphics Conference. 2013.
	
	\bibitem{unifiedcppshaders} Kerry A. Seitz Jr., Theresa Foley, Serban D. Porumbescu, John D. Owens. ``Supporting Unified Shader Specialization by Co-opting C++ Features''. arXiv:2109.14682 
	
	\bibitem{old} Frolov Vladimir, Sanzharov Vadim, Galaktionov Vladimir, Scherbakov Alexandr. An auto-programming approach to Vulkan. // Proceedings of the 31th International Conference on Computer Graphics and Machine Vision, CEUR Workshop Proceedings, Vol 3027, pp. 150-165. URL: http://ceur-ws.org/Vol-3027/paper14.pdf 
	
	\bibitem{taichi} ``Taichi Programming Language''. https://www.taichi-lang.org/
	
	\bibitem{numba} ``Numba''. JIT compiler. https://numba.readthedocs.io/en/stable/
	
	\bibitem{nvidiacompilers} Nvidia. ``NVIDIA HPC Fortran, C++ and C Compilers with OpenACC''.
	
	\bibitem{oneapi} J. Gold. ``OneAPI: Software Abstraction for a Heterog. Comp. World''. 
	
	\bibitem{drjit} Wenzel Jakob, Sebastien Speierer, Nicolas Roussel, and Delio Vicini. 2022. ``DR.JIT: a just-in-time compiler for differentiable rendering.'' ACM Trans. Graph. 41, 4, Article 124 (July 2022), 19 pages. https://doi.org/10.1145/3528223.3530099
	
	\bibitem{difftaichi} Yuanming Hu, Luke Anderson, Tzu-Mao Li, Qi Sun, Nathan Carr, Jonathan Ragan-Kelley, Fredo Durand. ``DiffTaichi: Differentiable Programming for Physical Simulation.'' ICLR 2020.
	
	\bibitem{diffhalide} Tzu-Mao Li, Michael Gharbi, Andrew Adams, Fredo Durand, and Jonathan Ragan-Kelley. ``Differentiable programming for image processing and deep learning in halide.'' ACM Trans. Graph. 37, 4, Article 139, 2018, 13 pages. https://doi.org/10.1145/3197517.3201383
	
	\bibitem{Reinhard05} Reinhard, E., Ward, G., Pattanaik, S., and Debevec, P. ``High Dynamic Range Imaging: Acquisition, Display and Image-Based Lighting'', 2005, chap. 6, pp. 187–221. Morgan Kaufmann Publishers Inc., first edition.
	
	\bibitem{Embree} Ingo Wald, Sven Woop, Carsten Benthin, Gregory S. Johnson, and Manfred Ernst. ``Embree: a kernel framework for efficient CPU ray tracing.'' ACM Trans. Graph. 33, 4, Article 143 (2014), 8 pages. https://doi.org/10.1145/2601097.2601199
	
	\bibitem{HydraRT} Ray Tracing Systems, Keldysh Institute of Applied Mathematics, Moscow State Uiversity. ``Hydra Renderer. Open source rendering system'', 2019.
	https://github.com/Ray-Tracing-Systems/HydraAPI
	
	\bibitem{kslicer} Frolov Vladimir, Sanzharov Vadim, Galaktionov Vladimir. kernel\_slicer. https://github.com/Ray-Tracing-Systems/kernel\_slicer
	
	\bibitem{vez} V-EZ, an open source, cross-platform wrapper, 2018. URL: https://github.com/GPUOpen-LibrariesAndSDKs/V-EZ
	
	\bibitem{vuh} Vuh. A Vulkan-based GPGPU computing framework, 2020 URL: https://github.com/Glavnokoman/vuh
	
	\bibitem{kompute} Kompute. The general purpose GPU compute framework for cross vendor graphics cards, 2021. URL: https://github.com/KomputeProject/kompute
	
	\bibitem{MLIR} Chris Lattner, Mehdi Amini, Uday Bondhugula, Albert Cohen, Andy Davis, Jacques Pienaar, River Riddle, Tatiana Shpeisman, Nicolas Vasilache, Oleksandr Zinenko. \textit{MLIR: A Compiler Infrastructure for the End of Moore's Law} // arXiv:2002.11054
	
	\bibitem{glslangValidator} KhronosGroup. glslangValidator // reference shader compiler for Vulkan. URL = https://github.com/KhronosGroup/glslang
	
	\bibitem{glslc} Google. Shaderc // A collection of tools, libraries and tests for shader compilation. URL = https://github.com/google/shaderc
	
	\bibitem{clvk} Kevin Petit. \textit{A prototype implementation of OpenCL 3.0 on top of Vulkan using clspv as the compiler.} // https://github.com/kpet/clvk
	
	\bibitem{Morton66} Morton G.M. \textit{A Computer Oriented Geodetic Data Base: and a New Technique in File Sequencing} // Research Report IBM Ltd., Ottawa, ON, Canada, 1966.
	
	\bibitem{Karras12} Tero Karras. Maximizing Parallelism in the Construction of BVHs, Octrees, and k-d Trees. // High Performance Graphics (2012)
	
	\bibitem{Karras13} Tero Karras and Timo Aila. \textit{Fast parallel construction of high-quality bounding volume hierarchies.} // In Proceedings of the 5th High-Performance Graphics Conference (HPG '13). 2013. Association for Computing Machinery, New York, NY, USA, 89–99. https://doi.org/10.1145/2492045.2492055
	
	\bibitem{fastbuild1} Christian Lauterbach, Michael Garland, Shubhabrata Sengupta et al. \textit{Fast BVH Construction on GPUs} // Computer Graphics Forum. — 2009.
	
	\bibitem{fastbuild2} Yan Gu, Yong He, Kayvon Fatahalian, Guy Blelloch. \textit{Efficient BVH Construction via Approximate Agglomerative Clustering} // Proceedings of the 5th High-Performance Graphics Conference. — HPG ’13. — Association for Computing Machinery, 2013. — P. 81–88. https://doi.org/10.1145/2492045.2492054.
	
	\bibitem{fastbuild3} Kirill Garanzha, Jacopo Pantaleoni, and David McAllister. \textit{Simpler and faster HLBVH with work queues.} // In Proceedings of the ACM SIGGRAPH Symposium on High Performance Graphics (HPG '11). 2011. Association for Computing Machinery, New York, NY, USA, 59–64. https://doi.org/10.1145/2018323.2018333
	
	\bibitem{fastbuild4} Chitalu, Francis M. \textit{Binary Ostensibly-Implicit Trees for Fast Collision Detection} // Computer Graphics Forum. — 2020. https://doi.org/10.1111/cgf.13948.
	
	\bibitem{fastbuild5} Ciprian Apetrei. \textit{Fast and Simple Agglomerative LBVH Construction} // EG UK Computer Graphics and Visual Computing (2014).
	
	\bibitem{MB22} Daniel Meister and Jiri Bittner. \textit{Performance Comparison of Bounding Volume Hierarchies for GPU Ray Tracing} // Journal of Computer Graphics Techniques (JCGT), vol. 11, no. 4, 1-19, 2022. Available online http://jcgt.org/published/0011/04/01/
	
	\bibitem{Aila2009} Timo Aila and Samuli Laine. \textit{Understanding the Efficiency of Ray Traversal on GPUs.} // Proc. High-Performance Graphics 2009.
	
	\bibitem{Reinhard05} Reinhard, E., Ward, G., Pattanaik, S., and Debevec, P. \textit{High Dynamic Range Imaging: Acquisition, Display and Image-Based Lighting}, 2005.
	
	\bibitem{Mantiuk08} Mantiuk R., Daly S., Kerofsky L. Display adaptive tone mapping //ACM SIGGRAPH 2008 papers. – 2008. – С. 1-10.
	
\end{thebibliography}

\end{document}